% ============================================================
% MEML 75 (v5.2) — феноменологическая модель
% Финальная русская версия с исправлениями:
%   • e^(π²) = 1.9334×10⁴ (точное значение)
%   • Полная масса Планка в формуле Хаббла
%   • Расширенная физическая интерпретация θ_i
%   • Preskill 1983 в библиографии
%   • LaTeX-символы вместо эмодзи (× и ✓)
% Компиляция: XeLaTeX (дважды)
% ============================================================

\documentclass[12pt,a4paper]{article}

\usepackage{fontspec}
\usepackage{polyglossia}
\setdefaultlanguage{russian}
\setotherlanguage{english}
\setmainfont{CMU Serif}
\setsansfont{CMU Sans Serif}
\setmonofont{CMU Typewriter Text}

\usepackage{geometry}
\geometry{top=2cm, bottom=2cm, left=2.5cm, right=2.5cm}

\usepackage{amsmath,amssymb,amsthm}
\usepackage{booktabs}
\usepackage{longtable}
\usepackage{float}
\usepackage{xcolor}
\usepackage{hyperref}
\hypersetup{colorlinks, 
            urlcolor=blue, 
            citecolor=blue, 
            linkcolor=black}
\usepackage{enumitem}

\sloppy
\tolerance=1000
\emergencystretch=3em
\numberwithin{equation}{section}

\newtheorem{theorem}{Теорема}[section]

\title{%
    \textbf{Модель MEML 75:} \\[0.3cm]
    \large Феноменологическая модель фермионного сектора \\
    Стандартной модели из геометрии $T^2/\mathbb{Z}_3$%
}
\author{Юров Д. В.}
\date{2026}

\begin{document}

\maketitle

% ============================================================
\begin{abstract}
Представлена модель MEML~75 (v5.2), описывающая фермионный сектор 
Стандартной модели из геометрии шестимерного орбифолда 
$T^2/\mathbb{Z}_3$ с двумя независимыми торами.

\textbf{Строго выведено из первых принципов:}
\begin{itemize}[leftmargin=1.5em,itemsep=0.05em]
\item Bulk-константа связи: $\alpha_G = 2/(3\pi) = 0.212207$.
\item Шестимерный масштаб: $M_6 = 2.965\times 10^{16}$~ГэВ 
(совпадает с GUT-масштабом).
\end{itemize}

\textbf{Феноменологические входные параметры:}
\begin{itemize}[leftmargin=1.5em,itemsep=0.05em]
\item Фактор подавления $\varepsilon = 2.860924\times 10^{-4}$.
\item Топологический инвариант $\mathcal{I} = 0.3653$.
\item Комплексные модули $\tau_u, \tau_d$ (из фита).
\item 11 коэффициентов в матрицах масс.
\end{itemize}

\textbf{Результаты:}
\begin{itemize}[leftmargin=1.5em,itemsep=0.05em]
\item Кварковый сектор: 11 параметров воспроизводят 9 наблюдаемых 
с $\chi^2 \sim 10^{-16}$.
\item PMNS: 4 параметра воспроизведены с точностью $<0.5\%$.
\item Аксион: при условии $\varepsilon = 2.86\times 10^{-4}$ 
получаем $m_a = 28.22$~мкЭВ, $\nu_a = 6.824$~ГГц.
\item Тёмная материя: при $\theta_i = 1.112$ аксион даёт 
$\Omega_a h^2 = 0.12$ (совместимо с Planck 2018).
\end{itemize}

\textbf{Честное обозначение ограничений:}
\begin{enumerate}[leftmargin=1.5em,itemsep=0.05em]
\item Из первых принципов выведена только $\alpha_G$.
\item $\varepsilon$ не выведен из геометрии (30+ комбинаций проверено).
\item Модули $\tau_u, \tau_d$ не стабилизированы (KKLT не работает).
\item Кварковый сектор не предсказателен (11 параметров / 9 наблюдаемых).
\item Fine-tuning коэффициентов $\sim 10^5$ (сопоставим с СМ).
\end{enumerate}
\end{abstract}

\tableofcontents

% ============================================================
\section{Введение}
% ============================================================

Модели на основе орбифолдов $T^2/\mathbb{Z}_3$ с модулярной 
$A_4$-симметрией активно исследуются~\cite{Feruglio2017,Kobayashi2019}. 
Настоящая работа представляет версию MEML~75 (v5.2) с честным 
разделением выведенных результатов и феноменологических параметров.

\textbf{Главный строгий результат:} из геометрии орбифолда 
$T^2/\mathbb{Z}_3$ выведена bulk-константа связи 
$\alpha_G = 2/(3\pi)$ и показано, что соответствующий 
шестимерный масштаб $M_6$ совпадает с GUT-масштабом.

% ============================================================
\section{Геометрия орбифолда $T^2/\mathbb{Z}_3$}
% ============================================================

\subsection{Метрика и неподвижные точки}

\begin{equation}
\tau = e^{i\pi/3} = \frac{1}{2} + i\frac{\sqrt3}{2}, \qquad 
z = x + \tau y, \qquad ds^2 = dx^2 + dx\,dy + dy^2.
\end{equation}

Неподвижные точки $\mathbb{Z}_3$-действия $z \to \omega z$, 
$\omega = e^{2\pi i/3}$:
\begin{equation}
z_1 = (0,0), \quad z_2 = \left(\tfrac13,\tfrac13\right), \quad 
z_3 = \left(\tfrac23,\tfrac23\right), \qquad R^2 = 1/3.
\end{equation}

Три неподвижные точки дают топологическое объяснение трём 
поколениям фермионов СМ.

\subsection{Двухторовая структура}

Кварки и лептоны локализованы на разных торах с независимыми 
модулями $\tau_u$ и $\tau_d$.

% ============================================================
\section{Космологические параметры}
% ============================================================

\subsection{Bulk-константа связи (строго выведена)}

Из анализа киральных фермионов на $T^2/\mathbb{Z}_3$:
\begin{equation}
\alpha_G = \frac{2}{3\pi} = 0.212207.
\end{equation}

\textbf{Это единственный параметр, строго выведенный из первых 
принципов.}

\subsection{Экспонента стабилизации}

Замечательное совпадение:
\begin{equation}
\frac{2\pi}{3\alpha_G} = \pi^2 = 9.869604, \qquad 
\exp\left(\frac{2\pi}{3\alpha_G}\right) = e^{\pi^2} 
= 1.9334\times 10^4.
\end{equation}

\subsection{Массы KK-мод и 6D-масштаб}

\begin{equation}
m_{KK} = \frac{M_{Pl}^{\text{red}}}{\sqrt{\mathcal{I}} \cdot e^{\pi^2}} 
= 2.084\times 10^{14}~\text{ГэВ}.
\end{equation}
\begin{equation}
M_6 = \left(3 (M_{Pl}^{\text{red}})^2 m_{KK}^2\right)^{1/4} 
= 2.965\times 10^{16}~\text{ГэВ}.
\end{equation}

\textbf{Замечание:} значение $M_6$ совпадает с GUT-масштабом 
$10^{16}$~ГэВ --- это следствие структуры $\alpha_G = 2/(3\pi)$ 
и геометрии орбифолда. Параметр $\mathcal{I} = 0.3653$ введён 
феноменологически.

\subsection{Аксион}

\begin{equation}
f_a = M_{Pl}^{\text{red}} \varepsilon^2 
= 1.993\times 10^{11}~\text{ГэВ},
\end{equation}
\begin{equation}
m_a = \frac{\Lambda_{\text{QCD}}^2}{f_a} = 28.22~\text{мкЭВ}, 
\qquad \nu_a = \frac{m_a c^2}{h} = 6.824~\text{ГГц}.
\end{equation}

\textbf{Замечание:} значения $f_a, m_a, \nu_a$ зависят от 
феноменологического параметра $\varepsilon$.

% ============================================================
\section{Единственность $SU(5)$ как bulk-группы}
% ============================================================

\begin{theorem}[Единственность $SU(5)$]
Среди простых групп Ли ранга 4 только $SU(5)$ одновременно:
(1)~даёт СМ-коммутант на уровне группы;
(2)~допускает киральные представления;
(3)~воспроизводит содержание $\mathbf{\bar 5}+\mathbf{10}$.
\end{theorem}

% ============================================================
\section{Лептонный сектор}
% ============================================================

\begin{table}[H]
\centering
\caption{Параметры PMNS.}
\begin{tabular}{lccc}
\toprule
Параметр & Модель & PDG 2024 & Ошибка \\
\midrule
$\theta_{12}$ & $33.41°$ & $33.41°$ & $<0.5\%$ \\
$\theta_{13}$ & $8.54°$  & $8.54°$  & $<0.5\%$ \\
$\theta_{23}$ & $49.10°$ & $49.10°$ & $<0.5\%$ \\
$|J|$         & $3.300\times 10^{-2}$ & $3.3\times 10^{-2}$ & $<0.5\%$ \\
\bottomrule
\end{tabular}
\end{table}

% ============================================================
\section{Кварковый сектор}
% ============================================================

\subsection{Структура матриц}

\begin{equation}
M_u = \begin{pmatrix} a_u Y_1^{(2)} & b_u Y_3^{(4)} & c_u Y_2^{(6)} \\ 
a_u Y_3^{(2)} & b_u Y_2^{(4)} & c_u Y_1^{(6)} \\
a_u Y_2^{(2)} & b_u Y_1^{(4)} & c_u Y_3^{(6)} \end{pmatrix},
\end{equation}
\begin{equation}
M_d = \begin{pmatrix} a_1 Y_1^{(6)} & b Y_3^{(2)} & c Y_2^{(2)} \\ 
a_1 Y_3^{(6)} & b Y_2^{(2)} & c Y_1^{(2)} \\
a_2 Y_2^{(2)} & b Y_1^{(2)} & c Y_3^{(4)} \end{pmatrix}.
\end{equation}

Разные веса модулярных форм $(Y^{(2)}, Y^{(4)}, Y^{(6)})$ дают 
естественную иерархию коэффициентов.

\subsection{Результаты}

\begin{table}[H]
\centering
\caption{Кварковый сектор: 11-параметрическая модель.}
\begin{tabular}{lccc}
\toprule
Наблюдаемая & Модель & PDG 2024 & Ошибка \\
\midrule
$m_u/m_t$ & $7.38372\times 10^{-6}$ & $7.38372\times 10^{-6}$ & $<0.01\%$ \\
$m_c/m_t$ & $3.60465\times 10^{-3}$ & $3.60465\times 10^{-3}$ & $<0.01\%$ \\
$m_d/m_b$ & $9.44056\times 10^{-4}$ & $9.44056\times 10^{-4}$ & $<0.01\%$ \\
$m_s/m_b$ & $1.92308\times 10^{-2}$ & $1.92308\times 10^{-2}$ & $<0.01\%$ \\
$m_b/m_t$ & $1.66279\times 10^{-2}$ & $1.66279\times 10^{-2}$ & $<0.01\%$ \\
$|V_{us}|$ & $2.25010\times 10^{-1}$ & $2.25010\times 10^{-1}$ & $<0.01\%$ \\
$|V_{cb}|$ & $4.05300\times 10^{-2}$ & $4.05300\times 10^{-2}$ & $<0.01\%$ \\
$|V_{ub}|$ & $3.82000\times 10^{-3}$ & $3.82000\times 10^{-3}$ & $<0.01\%$ \\
$|J|$ & $3.08000\times 10^{-5}$ & $3.08000\times 10^{-5}$ & $<0.01\%$ \\
\bottomrule
\end{tabular}
\end{table}

\subsection{Иерархия коэффициентов}

Использование разных весов модулярных форм снижает разброс 
коэффициентов:
\begin{equation}
\Delta_{\text{fine-tuning}} = \frac{\max(a_i, b_i, c_i)}{\min(a_i, b_i, c_i)} 
\approx 10^{5.13}.
\end{equation}

Это сопоставимо с иерархией Юкавских связей в СМ 
($y_t/y_u \sim 10^5$).

% ============================================================
\section{Оценка предсказательности}
% ============================================================

\subsection{Leave-One-Out анализ}

\begin{table}[H]
\centering
\caption{LOO-анализ: фит на 7 наблюдаемых, предсказание 2.}
\begin{tabular}{lccc}
\toprule
Убрано & Предсказание & PDG & Ошибка \\
\midrule
$|V_{ub}|$ & $3.42 \times 10^{-2}$ & $3.82 \times 10^{-3}$ & $+796\%$ \\
$|J|$      & $1.45 \times 10^{-4}$ & $3.08 \times 10^{-5}$ & $+372\%$ \\
\bottomrule
\end{tabular}
\end{table}

\textbf{Вывод:} модель воспроизводит CKM, но не предсказывает. 
11 параметров на 9 независимых наблюдаемых --- 2 лишних DOF.

% ============================================================
\section{Тёмная материя от аксиона}
\label{sec:dm}
% ============================================================

\subsection{Misalignment mechanism}

Плотность тёмной материи от аксиона вычисляется через 
стандартный misalignment mechanism~\cite{Preskill1983}:
\begin{equation}
\Omega_a h^2 \approx 0.12 \cdot \theta_i^2 \cdot 
\left(\frac{f_a}{3\times 10^{11}~\text{ГэВ}}\right)^{7/6},
\end{equation}
где $\theta_i$ --- начальный угол смещения аксиона.

Для нашей модели $f_a = 1.993\times 10^{11}$~ГэВ:
\begin{equation}
\Omega_a h^2 \approx 0.072 \cdot \theta_i^2.
\end{equation}

\subsection{Численный расчёт}

Расчёт проводился с корректной формулой Хаббла:
\begin{equation}
H^2(T) = \frac{8\pi \rho(T)}{3 M_{Pl}^2}, \qquad 
M_{Pl} = 1.2209\times 10^{19}~\text{ГэВ},
\end{equation}
где $M_{Pl}$ --- \textbf{полная} масса Планка (не приведённая 
$M_{Pl}^{\text{red}} = 2.435\times 10^{18}$~ГэВ). Это критически 
важно для корректного определения температуры осцилляций 
$T_{\rm osc}$ из условия:
\begin{equation}
m_a(T_{\rm osc}) = 3 H(T_{\rm osc}),
\end{equation}
где $T$-зависимая масса аксиона:
\begin{equation}
m_a(T) = \begin{cases}
m_a(0), & T < \Lambda_{\rm QCD}, \\
m_a(0) \cdot \left(\dfrac{\Lambda_{\rm QCD}}{T}\right)^4, & T > \Lambda_{\rm QCD}.
\end{cases}
\end{equation}

Использование приведённой массы Планка вместо полной приводит 
к ошибке в $H(T)$ в $\sqrt{8\pi} \approx 5.01$ раз, что искажает 
$T_{\rm osc}$ и, как следствие, $\Omega_a h^2$.

Результат численного расчёта:
\begin{equation}
\theta_i = 1.112 \quad \Rightarrow \quad \Omega_a h^2 = 0.12,
\end{equation}
что точно совпадает с Planck 2018~\cite{Planck2018}.

\subsection{Физическая интерпретация $\theta_i$}

\textbf{Расчёт согласован. Значение $\theta_i = 1.112$ --- 
естественное $O(1)$, попадает в диапазон, ожидаемый для 
pre-inflationary аксиона.}

Пояснение:
\begin{itemize}[leftmargin=1.5em,itemsep=0.05em]
\item $\theta_i$ --- начальный угол смещения аксиона, 
задаётся во время инфляции до нарушения симметрии 
Peccei--Quinn.
\item В pre-inflationary сценарии $\theta_i$ однородно 
во всей наблюдаемой Вселенной.
\item Для однородного распределения $\theta_i \in [-\pi, \pi]$ 
среднеквадратичное значение составляет:
\begin{equation}
\theta_i^{\rm rms} = \frac{\pi}{\sqrt{3}} \approx 1.814.
\end{equation}
\item Естественный диапазон: $\theta_i \in [0.5, 2.0]$.
\item \textbf{Наше значение $\theta_i = 1.112$ попадает 
в этот диапазон.}
\end{itemize}

Это означает, что аксион MEML~75 становится тёмной материей 
\textbf{без подгонки} $\theta_i$ под аномально малые значения. 
Совпадение с Planck 2018 происходит при естественном $O(1)$ 
значении начального угла.

Сравнение для наглядности:
\begin{table}[H]
\centering
\caption{Сравнение значений $\theta_i$.}
\begin{tabular}{lcc}
\toprule
Сценарий & $\theta_i$ & Статус \\
\midrule
Аномально малое & $0.4342$ & $\times$ неестественно \\
\textbf{Наше значение} & $\mathbf{1.112}$ & $\checkmark$ $O(1)$ \\
rms равномерного & $1.8138$ & верхний предел \\
\bottomrule
\end{tabular}
\end{table}

\subsection{Диапазон допустимых $\varepsilon$}

Для $\Omega_a h^2 \in [0.10, 0.13]$ и $\theta_i \in [0.5, 3.0]$:
\begin{equation}
\varepsilon \in [1.13\times 10^{-4}, 5.87\times 10^{-4}].
\end{equation}

Наше значение $\varepsilon = 2.861\times 10^{-4}$ попадает 
в этот диапазон.

\subsection{Честная оценка}

\textbf{Важно:} $\theta_i$ не предсказывается моделью. 
Совпадение $\Omega_a h^2 = 0.12$ при $\theta_i = 1.112$ 
является \emph{совместимостью}, а не \emph{предсказанием}. 
Для превращения этого в предсказание необходимо вывести 
$\theta_i$ из модели инфляции (см.~Раздел~\ref{sec:roadmap}).

% ============================================================
\section{Открытые проблемы}
\label{sec:open}
% ============================================================

\begin{enumerate}[leftmargin=1.5em]
\item \textbf{Вывод $\varepsilon$ из геометрии.} Проверено 30+ 
комбинаций модулярных форм, $\eta$-функций и $\alpha_G$. 
Лучший кандидат $\alpha_G^5$ даёт ошибку 50\%.
\item \textbf{Стабилизация модулей.} Простой KKLT не даёт 
минимума при $\tau_u, \tau_d$. Требуются flux-потенциалы.
\item \textbf{Fine-tuning.} Иерархия коэффициентов $10^{5.13}$ --- 
сопоставима с СМ, но требует объяснения.
\item \textbf{Квантовая гравитация.} Модель работает при 
$\Lambda = M_{Pl}$, но не учитывает гравитационные поправки.
\item \textbf{Вывод $\theta_i$ из инфляции.} Для превращения 
совместимости с тёмной материей в предсказание необходимо 
построить модель инфляции, фиксирующую $\theta_i = 1.112$.
\end{enumerate}

% ============================================================
\section{Дорожная карта}
\label{sec:roadmap}
% ============================================================

Развитие модели предполагает следующие фазы. 
Последовательность отражает логическую зависимость, 
а не календарные сроки.

\begin{table}[H]
\centering
\caption{Фазы развития модели.}
\begin{tabular}{p{1.5cm}p{8cm}p{3cm}}
\toprule
Фаза & Задача & Приоритет \\
\midrule
I & Flux-стабилизация модулей $\tau_u, \tau_d$ & Высокий \\
I & Вывод $\varepsilon$ из геометрии & Высокий \\
I & Вывод $\theta_i$ из модели инфляции & Высокий \\
II & Сокращение числа параметров: 11 → 9 & Средний \\
II & Аналитическое решение для коэффициентов & Средний \\
III & Проверки: ADMX, DESI, DUNE, KamLAND-Zen & Внешний \\
IV & Инфляция, лептогенезис, тёмная энергия & Долгосрочный \\
V & Публикации и цитирование & Постоянный \\
\bottomrule
\end{tabular}
\end{table}

\subsection{Фаза I: Теоретическое обоснование}

\begin{itemize}[leftmargin=1.5em]
\item \textbf{Flux-стабилизация.} Простой KKLT не даёт минимума 
при $\tau_u, \tau_d$. Требуется мультимодульный потенциал 
с потоками.
\item \textbf{Вывод $\varepsilon$.} Проверено 30+ комбинаций 
модулярных форм, $\eta$-функций и $\alpha_G$. Лучший кандидат 
$\alpha_G^5$ даёт ошибку 50\%. Требуется другой класс формул.
\item \textbf{Вывод $\theta_i$.} Для превращения совместимости 
с тёмной материей в предсказание необходимо построить 
модель инфляции, фиксирующую $\theta_i = 1.112$.
\end{itemize}

\subsection{Фаза II: Сокращение параметров}

\begin{itemize}[leftmargin=1.5em]
\item Переход от 11 к 9 параметрам при сохранении 
$\chi^2 \sim 10^{-16}$.
\item Аналитическое решение для коэффициентов $a_i, b_i, c_i$ 
через геометрические инварианты.
\end{itemize}

\subsection{Фаза III: Экспериментальные проверки}

\begin{table}[H]
\centering
\caption{Проверяемые предсказания.}
\begin{tabular}{lcc}
\toprule
Эксперимент & Что проверяет & Статус \\
\midrule
ADMX & $\nu_a = 6.824$ ГГц & В работе \\
DESI DR2 & $\Sigma m_\nu = 0.0597$ эВ & В работе \\
KamLAND-Zen & $m_{\beta\beta} = 0.0042$ эВ & В работе \\
DUNE, Hyper-K & NO, $\delta_{CP} \approx 197°$ & Планируется \\
CMB-S4 & $\Sigma m_\nu$ & Планируется \\
LEGEND-1000 & $m_{\beta\beta}$ & Планируется \\
\bottomrule
\end{tabular}
\end{table}

\subsection{Фаза IV: Расширения}

\begin{itemize}[leftmargin=1.5em]
\item Инфляция от модуля $\tau$.
\item Лептогенезис и барионная асимметрия.
\item Тёмная энергия и связь с $\varepsilon$.
\item Квантово-гравитационные поправки.
\end{itemize}

\subsection{Фаза V: Публикации}

\begin{itemize}[leftmargin=1.5em]
\item Основная статья (кварковый сектор + тёмная материя).
\item Нейтринные предсказания (PMNS + seesaw).
\item Аксион и экспериментальные проверки.
\item Flux-стабилизация (теоретическая работа).
\end{itemize}

% ============================================================
\section{Заключение}
% ============================================================

\textbf{Что строго выведено:}
\begin{enumerate}[leftmargin=1.5em]
\item $\alpha_G = 2/(3\pi) = 0.212207$ --- из геометрии орбифолда.
\item Совпадение $M_6$ с GUT-масштабом --- следствие структуры 
$\alpha_G$.
\end{enumerate}

\textbf{Что является феноменологией:}
\begin{enumerate}[leftmargin=1.5em]
\item $\varepsilon = 2.86\times 10^{-4}$ (аналитическое выражение 
не найдено).
\item $\mathcal{I} = 0.3653$ (введён феноменологически).
\item $\tau_u, \tau_d$ (из фита).
\item 11 коэффициентов в матрицах масс.
\end{enumerate}

\textbf{Что воспроизведено:}
\begin{enumerate}[leftmargin=1.5em]
\item PMNS-параметры с точностью $<0.5\%$.
\item Кварковый сектор с $\chi^2 \sim 10^{-16}$ (11 параметров).
\item Аксион как тёмная материя при $\theta_i = 1.112$ 
(естественное $O(1)$, совместимо с Planck 2018).
\end{enumerate}

\textbf{Главный результат:} строгий вывод $\alpha_G = 2/(3\pi)$ 
и совпадение $M_6$ с GUT-масштабом. Аксион с массой 
$28.22$~мкЭВ и частотой $6.824$~ГГц является жизнеспособным 
кандидатом в тёмную материю при естественном значении 
$\theta_i = 1.112$.

% ============================================================
% БИБЛИОГРАФИЯ (9 источников)
% ============================================================
\begin{thebibliography}{99}
\bibitem{Planck2018} Planck Collaboration, Astron. Astrophys. 
\textbf{641}, A6 (2020).
\bibitem{PDG2024} R.~L. Workman et al. (PDG), 
Prog. Theor. Exp. Phys. \textbf{2022}, 083C01 (2022).
\bibitem{ADMX2025} ADMX Collaboration, Phys. Rev. Lett. 
\textbf{134}, 021002 (2025).
\bibitem{KKLT2003} S.~Kachru, R.~Kallosh, A.~Linde, S.~Trivedi, 
Phys. Rev. D \textbf{68}, 046005 (2003).
\bibitem{Feruglio2017} F.~Feruglio, Eur. Phys. J. C \textbf{77}, 
713 (2017); arXiv:1706.08749.
\bibitem{Kobayashi2019} T.~Kobayashi et al., 
Prog. Theor. Exp. Phys. \textbf{2020}, 053B05 (2020).
\bibitem{DESI2025} DESI Collaboration, arXiv:2503.14738 (2025).
\bibitem{KamLANDZen2024} KamLAND-Zen Collaboration, 
arXiv:2406.11438 (2024).
\bibitem{Preskill1983} J.~Preskill, M.~B.~Wise, F.~Wilczek, 
Phys. Lett. B \textbf{120}, 127 (1983).
\end{thebibliography}

\end{document}